\scanpage{99}
\mspage{72}
\[
(\widetilde F_x(u,v))(a_1,\ldots,a_p,b_1,\ldots,b_p)=(F_x(u-v))(a_1-u,\ldots,a_p-u,b_1-v,\ldots,b_p-v)
\]
where $a_i,b_i\in T_x$. We denote by $\mathcal H_x^p$ the set of all operators on $T_x$ which can be obtained from the mapping $F_x$ as above. Now note that $\mathbb H_xJ$ is the matrix whose elements are the homogeneous parts of the elements of the matrix corresponding to $J$. Hence if we denote by $\mathcal F_x^p$ the ideal of $\mathcal G^p$ consisting of all $J\in\mathcal G^p$ with $F_x=0$ and $\xi(x)=0$ then we obtain

\uline{Theorem 4.6.6.} $\mathcal H_x^p$ is an algebra under composition of operators. The map $\mathbb H_x:\mathcal G^p\longrightarrow\mathcal H_x^p$ is a homomorphism of algebras with kernel $\mathcal F_x^p$.

Since the space $T_x$ has an inner product we have an isomorphism $T_x\longrightarrow T_x^*$ which induces a natural isomorphism between $A_x^p\otimes A_x^p$ and the space of linear transformations $A_x^p\longrightarrow A_x^p$. This isomorphism provides $A_x^p\otimes A_x^p$ with an algebraic structure. Thus $\mathcal S_x^p$, the set of $C^\infty$ maps $S_x\longrightarrow A_x^p\otimes A_x^p$, is an algebra. Note that $\mathcal S_x^p$ can be represented as the algebra of $\binom{m}{p}\mathbin{\mathrm x}\binom{m}{p}$ matrices of $C^\infty$ functions on $S_x$.

Consider the map $\mathbb S_x:\mathcal H_x^p\longrightarrow\mathcal S_x^p$ defined by
\[
\tag{4.6.7}
(\mathbb S_x\mathbb H_xJ)(u)=\xi(x)+\int_{T_x}e^{iu\cdot v}\frac{F_x(v)}{\lVert v\rVert^m}\,dV_v
\]
\mspage{98}
